\section{EXPONENTIALS AND ITERATED LOGARITHMS}

We have already (V, p.~229) defined the iterated logarithms \(l_n(x)\) by the conditions \(l_0(x) = x\) , \(l_n(x) = \log(l_{n-1}(x))\) for \(n \geqslant 1\) . In the same way one defines the iterated exponentials \(e_n(x)\) by the conditions \(e_0(x) = x\) , \(e_n(x) = \exp(e_{n-1}(x))\) for \(n \geqslant 1\) . It is immediate, by induction on \(n\) , that \(l_n(x)\) is the inverse function of \(e_n(x)\) , defined for \(x > e_{n-1}(0)\) , and that \(e_m(e_n(x)) = e_{m+n}(x)\) , \(l_m(l_n(x)) = l_{m+n}(x)\) . By the relations \(\log x \ll x^\mu \ll e^x\) for every \(\mu > 0\) , one has, for \(n \geqslant 1\) ,

\[
l _ {n} (x) \prec (l _ {n - 1} (x)) ^ {\mu} \quad \text {   for   every   } \mu > 0\tag{5}
\]

\[
\begin{array}{c} e _ {n - 1} (x ^ {1 + \beta}) \prec \prec e _ {n} (x ^ {1 + \delta}) \prec \prec e _ {n} ((1 - \gamma) x) \prec \prec (e _ {n} (x)) ^ {\mu} \\ \prec \prec e _ {n} ((1 + \alpha) x) \prec \prec e _ {n} (x ^ {1 + \beta}) \end{array}\tag{6}
\]

for all \(\mu > 0, \alpha > 0, \beta > 0, 0 < \gamma < 1, 0 < \delta < 1\) , these numbers being otherwise arbitrary (cf.~V, p.~218, prop. 11).

We have already seen (V, p.~229) that, for \(n \geqslant 1\) , one has

\[
{\frac {d}{d x}} \big (l _ {n} (x) \big) = \prod_ {i = 0} ^ {n - 1} {\frac {1}{l _ {i} (x)}}.\tag{7}
\]

Similarly for \(n\geqslant 1\)

\[
\frac {d}{d x} \big (e _ {n} (x) \big) = \prod_ {i = 1} ^ {n} e _ {i} (x).\tag{8}
\]

whence, by prop. 8 of V, p.~233, for every \(\mu > 0\)

\[
\int_ {a} ^ {x} e _ {n} (t ^ {\mu}) d t \sim \frac {x}{\mu} e _ {n} (x ^ {\mu}) \prod_ {i = 0} ^ {n - 1} \frac {1}{e _ {i} (x ^ {\mu})}\tag{9}
\]

\[
\int_ {x} ^ {+ \infty} \frac {d t}{e _ {n} (t ^ {\mu})} \sim \frac {x}{\mu} \prod_ {i = 0} ^ {n} \frac {1}{e _ {i} (x ^ {\mu})}.\tag{10}
\]

One can show that if \(f\) is any (H) function such that \(f \gg 1\) then there exist two integers \(m\) and \(n\) such that

\[
l _ {m} (x) \prec f (x) \prec e _ {n} (x)
\]

(V, p.~263, exerc. 1 and p.~264, exerc. 5) On the other hand, one can define increasing functions \(g(x)\) (which are no longer (H) functions) such that \(g(x) \gg e_n(x)\) for every \(n > 0\) , or \(1 \ll g(x) \ll l_m(x)\) for every \(m > 0\) (V, p.~265, exerc. 8, 9 and 10).

With the help of the iterated logarithms we shall show that one can define a comparison scale \(\mathcal{E}\) (for \(x\) tending to \(+\infty\) ) of (H) functions, which are \(>0\) on a neighbourhood of \(+\infty\) and satisfy the following conditions:

\begin{enumerate}
\def\labelenumi{\alph{enumi})}
\item
  the product of any two functions in E belongs to E;
\item
  \(f^{\mu} \in \mathcal{E}\) for every function \(f \in \mathcal{E}\) and every real number \(\mu\) ;
\item
  for every function \(f \in \mathcal{E}\) , \(\log f\) is a linear combination of a finite number of functions in \(\mathcal{E}\) ;
\end{enumerate}

\(d\) ) for every function \(f \in \mathcal{E}\) , apart from the constant 1, \(e^f\) is equivalent to a function in \(\mathcal{E}\) .

First we consider the set \(E_{0}\) of functions of the form \(\prod_{m=0}^{\infty}\left(l_{m}(x)\right)^{\alpha_{m}}\) , where the \(\alpha_{m}\) are real numbers, zero apart from for a finite number of indices m; it is immediate, from (5) (V, p.~253) that these functions form a comparison scale which satisfies conditions a), b) and c). Now we define, by recursion on n, the set \(E_{n}\) (for \(n \geqslant 1\) ) formed by the constant 1 and by the functions of the form \(\exp\left(\sum_{k=1}^{p}a_{k}f_{k}\right)\) , where p is an arbitrary integer \textgreater{} 0, the functions \(f_{k}\) ( \(1 \leqslant k \leqslant p\) ) are functions in \(E_{n-1}\) such that \(f_{1} \gg f_{2} \gg \cdots \gg f_{p} \gg 1\) , and the \(a_{k}\) are real numbers \(\neq 0\) ; we show by induction that \(E_{n}\) is a comparison scale satisfying a), b) and c) and containing \(E_{n-1}\) . In the first place, the relation \(E_{n-1} \subset E_{n}\) holds for n = 1, since the logarithm of any nonconstant function in \(E_{0}\) is of the form \(\sum_{k=1}^{p}a_{k}f_{k}\) , where the \(f_{k}\) are iterated logarithms, and so \(\gg 1\) ; on the other hand, if \(E_{n-2} \subset E_{n-1}\) one deduces from the definition of \(E_{n}\) that \(E_{n-1} \subset E_{n}\) ; this definition furthermore shows that \(E_{n}\) satisfies a), b) and c). It remains to see that \(E_{n}\) is a comparison scale: since the quotient of two functions in \(E_{n}\) again belongs to \(E_{n}\) it suffices to prove that every function f of \(E_{n}\) , apart from the constant 1, cannot be equivalent to a constant \(\neq 0\) . Now one has \(\log f = \sum_{k=1}^{p}a_{k}f_{k} \sim a_{1}f_{1}\) by construction, and since \(f_{1} \gg 1\) , log f tends to \(\pm\infty\) , so f tends to 0 or to \(+\infty\) as x tends to \(+\infty\) .

This being so, if E is the union of the \(E_{n}\) for \(n \geqslant 0\) , then E is a comparison scale, for two functions in E belong to the same scale \(E_{n}\) ; for the same reason, E satisfies a), and it is clear that it also satisfies b) and c). Finally, if \(f \in E\) there exists an n such that \(f \in E_{n}\) ; if f is not the constant 1 then \(f(x)\) tends to 0 or to \(+\infty\) as x tends to \(+\infty\) ; in the first case \(e^{f} \sim 1\) and in the second, \(e^{f}\) belongs to \(E_{n+1}\) by definition, and so to E.

Remark. Despite the practical usefulness of the scale E which we have just defined, it is easy to give examples of (H) functions which have no principal part with respect to E. Indeed, if f is an (H) function such that \(f \sim ag\) , where a is a constant \textgreater{} 0 and \(g \in E\) , then \(\log f - \log g - \log a\) tends to 0 with 1/x, so \(\log f\) admits, relative to E, an asymptotic expansion whose remainder tends to 0, by property c). Now, if one considers for example the (H) function \(f(x) = e_{2}\left(x + \frac{1}{x}\right)\) one has \(\log f(x) = \exp\left(x + \frac{1}{x}\right)\) , so the asymptotic expansions of \(\log f\) relative to E are of the form

\[
\log f (x) = e ^ {x} + \frac {e ^ {x}}{x} + \frac {1}{2 !} \frac {e ^ {x}}{x ^ {2}} + \dots + \frac {1}{n !} \frac {e ^ {x}}{x ^ {n}} + o \left(\frac {e ^ {x}}{x ^ {n}}\right) \quad (n \text {   an   integer   } > 0).
\]

It is clear that the remainder in this expansion is equivalent to \(\frac{1}{(n+1)!}\frac{e^{x}}{x^{n+1}}\) , so does not tend to 0. Hence f does not have a principal part relative to E.

